\documentclass[10pt]{article}
\input{common}
\title{RevIN Experiment 2: Normalized Backpropagation}
\author{}
\date{}
\begin{document}\maketitle
\section{Theoretical overview and goal}
With instance statistics $(\mu_x,\sigma_x)$, data-space MSE is
$\lVert\hat y-y\rVert^2$, whereas normalized backpropagation minimizes
\[
 \operatorname{nMSE}(\hat y,y\mid x)=
 H^{-1}\sum_{h=1}^{H}
 \left(\frac{\hat y_h-y_h}{\sigma_x+\varepsilon}\right)^2.
\]
The latter gives comparable gradient weight to low- and high-scale windows.
Predictions are still denormalized and evaluated with ordinary data-space MSE,
so only optimization weighting changes.

\goal{Test whether scale-balanced gradients improve generalization independently
of the forward normalization layer, and whether RevIN's affine parameters alter
that effect.}

\section{Implementation details}
Eight paired contrasts are run:
\begin{center}
\begin{tabular}{lll}
\toprule
Forward transform & Data-space loss & Normalized loss \\
\midrule
None & \code{none\_mse} & \code{none\_nmse} \\
Global standardization & \code{standard\_mse} & \code{standard\_nmse} \\
Mean only & \code{mean\_only\_mse} & \code{mean\_only\_nmse} \\
Scale only & \code{scale\_only\_mse} & \code{scale\_only\_nmse} \\
Non-affine instance & \code{instance\_mse} & \code{instance\_nmse} \\
Median--MAD & \code{median\_mad\_mse} & \code{median\_mad\_nmse} \\
Affine RevIN & \code{revin\_mse} & \code{revin\_nmse} \\
Global MIN & \code{min\_mse} & \code{min\_nmse} \\
\bottomrule
\end{tabular}
\end{center}
Architecture, seed, batch composition, learning rate, and epoch budget are held
fixed within every pair. The paper comparison focuses on MSE versus nMSE.

The loader merges shared, RevIN-scoped, and run-specific user exclusions and
records them in every seed directory. This does not replace the separate
constant-window audit required by nMSE.

\textbf{Loss implementation:} \code{src/utils/pipeline.py}\\
\textbf{Experiment orchestration:} \code{src/scripts/experiment.py}

The root \code{revin.slurm} front runs signature-aware training and table
stages; resubmission skips current seed and table artifacts.

Small covers five datasets, the three shared default settings, PatchTST, and 16
component methods; full adds all four original ETT panels and two settings,
for nine datasets and five settings. Ultra adds DLinear. Publication runs
use three seeds and 10,000 steps, validating every 1,000 steps. Tables report
data-space test MSE and seed deviations. Validation selection and the optimistic
test oracle use every scheduled method by default.

\section{Interpretation}
nMSE can improve data-space MSE by balancing heterogeneous scales. MSE and nMSE
training curves have different units; select with a common validation metric.
\end{document}
